\SourceLeaf{252}%
\(\mathrm{d}p,\mathrm{d}q,\mathrm{d}p',\mathrm{d}q'\) will be
\[
\BareRoot{E}\mathord{.}\mathrm{d}p,\qquad\BareRoot{G}\mathord{.}\mathrm{d}q,\qquad\BareRoot{E'}\mathord{.}\mathrm{d}q',\qquad\BareRoot{G'}\mathord{.}\mathrm{d}q'
\]
We shall denote by \(M,N,M',N'\) the angles that the directions of these elements make with an arbitrary fixed direction, measuring in the sense in which the second lies with respect to the first, so that \(\sin(N-M)\) is positive. We suppose, as we may, that the fourth lies in the same sense with respect to the third, so that \(\sin(N'-M')\) is positive as well. With these conventions, consider another point at an infinitesimal distance from the first, with variable values
\[
p+\mathrm{d}p,\quad q+\mathrm{d}q,\quad p'+\mathrm{d}p',\quad q'+\mathrm{d}q'
\]
A little attention shows that the following holds in general, independently of the values of the increments \(\mathrm{d}p,\mathrm{d}q,\mathrm{d}p',\mathrm{d}q'\),
\[
\begin{aligned}
&\BareRoot{E}\mathord{.}\mathrm{d}p\mathord{.}\sin M+\BareRoot{G}\mathord{.}\mathrm{d}q\mathord{.}\sin N\\
&=\BareRoot{E'}\mathord{.}\mathrm{d}p'\mathord{.}\sin M'+\BareRoot{G'}\mathord{.}\mathrm{d}q'\mathord{.}\sin N'
\end{aligned}
\]
since each expression is simply the distance of the new point from the line from which the direction angles are measured. By the notation already introduced above we have \(N-M=\omega\). By analogy we set \(N'-M'=\omega'\), and in addition \(N-M'=\psi\). The equation just found can therefore be written in the following form
\[
\begin{aligned}
&\BareRoot{E}\mathord{.}\mathrm{d}p\mathord{.}\sin(M'-\omega+\psi)+\BareRoot{G}\mathord{.}\mathrm{d}q\mathord{.}\sin(M'+\psi)\\
&=\BareRoot{E'}\mathord{.}\mathrm{d}p'\mathord{.}\sin M'+\BareRoot{G'}\mathord{.}\mathrm{d}q'\mathord{.}\sin(M'+\omega')
\end{aligned}
\]
or in this form
\[
\begin{aligned}
&\BareRoot{E}\mathord{.}\mathrm{d}p\mathord{.}\sin(N'-\omega-\omega'+\psi)+\BareRoot{G}\mathord{.}\mathrm{d}q\mathord{.}\sin(N'-\omega'+\psi)\\
&=\BareRoot{E'}\mathord{.}\mathrm{d}p'\mathord{.}\sin(N'-\omega')+\BareRoot{G'}\mathord{.}\mathrm{d}q'\mathord{.}\sin N'
\end{aligned}
\]

Since the equation clearly must be independent of the initial direction, we may choose that direction at will. Setting \(N'=0\) in the second form or \(M'=0\) in the first, we obtain the following equations:
\[
\begin{aligned}
\BareRoot{E'}\mathord{.}\sin\omega'\mathord{.}\mathrm{d}p'&=\BareRoot{E}\mathord{.}\sin(\omega+\omega'-\psi)\mathord{.}\mathrm{d}p+\BareRoot{G}\mathord{.}\sin(\omega'-\psi)\mathord{.}\mathrm{d}q\\
\BareRoot{G'}\mathord{.}\sin\omega'\mathord{.}\mathrm{d}q'&=\BareRoot{E}\mathord{.}\sin(\psi-\omega)\mathord{.}\mathrm{d}p+\BareRoot{G}\mathord{.}\sin\psi\mathord{.}\mathrm{d}q
\end{aligned}
\]
Since these equations must be identical with
\[
\begin{aligned}
\mathrm{d}p'&=\alpha\mathrm{d}p+\beta\mathrm{d}q\\
\mathrm{d}q'&=\gamma\mathrm{d}p+\delta\mathrm{d}q
\end{aligned}
\]
